\capitulo{1}{Introducción}

El sueño cobra un papel fundamental en la salud humana, influyendo directamente en las funciones fisiológicas, cognitivas y emocionales. La existencia de trastornos del sueño, como el insomnio o la apnea obstructiva del sueño, se ha asociado con un mayor riesgo de desarrollar enfermedades cardiovasculares, metabólicas y mentales, así como con una disminución significativa en la calidad de vida de quienes los padecen.

Se estima que hasta un 30\% de la población mundial presenta algún tipo de alteración en el sueño. Aunque no siempre se consideran enfermedades graves, estos trastornos constituyen un problema de salud pública debido a su alta prevalencia, su impacto funcional y las comorbilidades que pueden desencadenar. Además, conllevan importantes costes económicos derivados de las bajas laborales, la disminución del rendimiento y el aumento del riesgo de accidentes de tráfico o laborales.


Actualmente, el diagnóstico de estos trastornos requiere herramientas clínicas especializadas y costosas, siendo la polisomnografía el estándar de referencia para la detección de alteraciones respiratorias o neurológicas durante el sueño. En otros casos, el diagnóstico se basa en cuestionarios subjetivos, cuya fiabilidad depende en gran medida de la percepción del paciente. En este contexto, los avances en aprendizaje automático han abierto nuevas oportunidades para el cribado y la detección temprana de alteraciones del sueño mediante el uso de variables fácilmente recopilables como la edad, la ocupación, los hábitos de sueño o la presión arterial.


Este Trabajo de Fin de Grado plantea el desarrollo de un sistema predictivo basado en aprendizaje automático, capaz de estimar la probabilidad de que un individuo presente un trastorno del sueño a partir de variables clínicas y de estilo de vida. Para ello, se ha utilizado un conjunto de datos público, se han evaluado distintos modelos de clasificación y se ha implementado una aplicación interactiva que permite introducir información del usuario y obtener una predicción explicada mediante técnicas de interpretabilidad.

El sistema propuesto pretende servir como herramienta de cribado accesible y explicable, aportando valor tanto a profesionales sanitarios como a usuarios particulares, especialmente en contextos de atención primaria o educación para la salud. Gracias a su diseño accesible, su carácter interpretable y su automatización, puede facilitar una primera valoración objetiva sin necesidad de equipamiento médico especializado.

Este documento recoge todas las fases del proyecto: desde el análisis de datos y la selección del modelo hasta la implementación de la aplicación web. Finalmente, se analizan los resultados obtenidos, las limitaciones detectadas y las posibles líneas de mejora para futuros desarrollos.